% ============================================================
%  قانون مدیریت دوران گذار ایران (طرح جامع مجلس مهستان)
%  نسخه‌ی نهایی - منطبق بر تمامی بندها و پیوست‌های سند اصلی
%  XePersian Book — Compiled with XeLaTeX
% ============================================================

\documentclass[12pt,a4paper,oneside]{report}

\usepackage{geometry}
\geometry{a4paper, top=2.5cm, bottom=2.5cm, left=2.5cm, right=2.5cm}

\usepackage[most]{tcolorbox}
\usepackage{xcolor}
\usepackage{titlesec}
\usepackage{hyperref}
\usepackage{enumitem}
\usepackage{fancyhdr}
\usepackage{longtable}
\usepackage{booktabs}

% ── Color Palette ───────────────────────────────────────────
\definecolor{MahestanBlue}{RGB}{21,101,192}
\definecolor{MahestanDarkBlue}{RGB}{13,71,161}
\definecolor{MahestanGold}{RGB}{249,168,37}
\definecolor{MahestanGray}{RGB}{97,97,97}

\hypersetup{
    colorlinks=true,
    linkcolor=MahestanBlue,
    urlcolor=MahestanBlue,
    pdftitle={قانون مدیریت دوران گذار ایران - طرح جامع مهستان}
}

% ── XePersian ────────────────────────────────────────────────
\usepackage{xepersian}
\settextfont[Scale=1.1]{Vazirmatn}
\setdigitfont{Vazirmatn}

% ── Styles ──────────────────────────────────────────────────
\newtcolorbox{articlebox}[2][]{%
  enhanced, parbox=false,
  before skip=10pt, after skip=10pt,
  colback=white, colframe=MahestanBlue!30,
  arc=2pt, boxrule=0.5pt,
  borderline west={3pt}{0pt}{MahestanBlue},
  left=10pt, right=10pt, top=8pt, bottom=8pt,
  title={\rl{\textbf{#2}}},
  fonttitle=\bfseries\small\color{MahestanDarkBlue},
  attach boxed title to top right={yshift=-2mm, xshift=-5mm},
  boxed title style={colback=MahestanBlue!10, colframe=MahestanBlue!30, arc=2pt},
  #1
}

\titleformat{\chapter}[display]
  {\normalfont\bfseries\Huge\color{MahestanDarkBlue}\raggedleft}
  {\small\color{MahestanGray} فصل \thechapter}
  {0.5ex}
  {\titlerule\vspace{1ex}}
  [\vspace{1ex}\titlerule]

\pagestyle{fancy}
\fancyhf{}
\fancyhead[R]{\small\color{MahestanGray} قانون جامع مدیریت دوران گذار (طرح مهستان)}
\fancyfoot[C]{\thepage}
\renewcommand{\headrulewidth}{0.4pt}

% ── Content Starts Here ──────────────────────────────────────
\begin{document}

\begin{titlepage}
    \centering
    \vspace*{2cm}
    {\Huge\bfseries\color{MahestanDarkBlue} قانون جامع مدیریت دوران گذار ایران \par}
    \vspace{1cm}
    {\huge\bfseries\color{MahestanGold} طرح مجلس مهستان \par}
    \vspace{0.8cm}
    {\Large (تدوین بر اساس تمامی بندها و پیوست‌های نهادی) \par}
    \vfill
    \begin{tcolorbox}[colback=MahestanBlue!5, colframe=MahestanBlue, arc=10pt, width=0.8\textwidth]
        \centering\itshape
        این قانون مجموعه‌ای فراگیر از تمامی ارکان، وظایف و محدودیت‌های نهادهای دوران گذار است که بر مبنای بیانیه‌ی چارچوب نهادی گذار تدوین شده است.
    \end{tcolorbox}
    \vfill
    {\small اسفند ۱۴۰۴ / مارس ۲۰۲۶ \par}
\end{titlepage}

\tableofcontents
\newpage

\chapter{اصول کلی و مبانی حاکمیت}

\begin{articlebox}{ماده ۱: ماهیت و بازه‌ی زمانی دوران گذار}
گذار ایران یک فرایند موقتی است که با سقوط نظام پیشین آغاز شده و حداکثر ظرف ۲۴ ماه با استقرار نهادهای ناشی از قانون اساسی نهایی خاتمه می‌یابد. تمدید این دوره تنها برای یک بار و به مدت ۶ ماه با رأی ۳/۴ مجلس مهستان میسر است.
\end{articlebox}

\begin{articlebox}{ماده ۲: منبع حاکمیت و مشروعیت}
حاکمیت ملی منحصراً متعلق به مردم ایران است. مشروعیت تمامی نهادهای دوران گذار ناشی از تأیید عمومی، انتخابات آزاد و یا توافق اجماعی جریان‌های ملی در فاز صفر است.
\end{articlebox}

\begin{articlebox}{ماده ۳: اصول بنیادین غیرقابل‌نقض (اصول ۱۴گانه)}
تمامی نهادهای دوران گذار و قانون اساسی نهایی ملزم به رعایت اصول ذیل هستند که تحت هیچ شرایطی قابل تعلیق نمی‌باشند:
\begin{enumerate}
    \item حاکمیت مردم بر اساس انتخابات آزاد.
    \item جدایی نهاد دین از نهاد دولت (بی‌طرفی دولت در قبال عقاید).
    \item برابری تمامی شهروندان بدون تبعیض جنسیتی، قومی، مذهبی یا گرایش جنسی.
    \item آزادی‌های بنیادین (بیان، مطبوعات، تجمع، تشکل و عقیده).
    \item ممنوعیت مطلق شکنجه و مجازات‌های خشن و غیرانسانی.
    \item حق دادرسی عادلانه، فرض بی‌گناهی و ممنوعیت بازداشت بدون حکم قضایی.
    \item حقوق اقلیت‌های قومی و زبانی و تضمین آموزش به زبان مادری.
    \item برابری کامل حقوقی، اجتماعی و اقتصادی زنان و مردان.
    \item استقلال کامل قوه قضاییه.
    \item حفظ تمامیت ارضی ایران بر پایه حکمرانی دموکراتیک و رفراندوم.
    \item حاکمیت قانون بر تمامی افراد و نهادها.
    \item ارجحیت عدالت و دادخواهی بر انتقام‌جویی شخصی یا گروهی.
    \item آزادی مطلق اطلاعات و ممنوعیت هرگونه سانسور و فیلترینگ.
    \item تعهد به صلح و حل مسالمت‌آمیز اختلافات بین‌المللی.
\end{enumerate}
\end{articlebox}

\chapter{شورای هماهنگی گذار (فاز صفر و خلأ اولیه)}

\begin{articlebox}{ماده ۴: مدیریت خلأ قدرت (روز صفر تا هفته ۵)}
جهت جلوگیری از هرج‌ومرج در فاصله‌ی سقوط نظام تا تشکیل مجلس مهستان، «شورای هماهنگی گذار» با مأموریت محدود تشکیل می‌شود.
\end{articlebox}

\begin{articlebox}{ماده ۵: ترکیب و مأموریت شورا}
۱. شورا متشکل از ۱۵ تا ۲۰ نفر از چهره‌های منتخب ائتلاف‌های اپوزیسیون با رعایت تنوع قومی و جنسیتی است.
۲. مأموریت شورا منحصراً شامل: برگزاری انتخابات مجلس مهستان، صدور ۱۰ فرمان اضطراری، تثبیت نظم عمومی و حفاظت از اسناد ملی است.
۳. شورا حق تصویب قوانین ساختاری، انتصابات بلندمدت یا امضای معاهدات بین‌المللی را ندارد.
۴. شورا بلافاصله پس از تشکیل مجلس مهستان منحل می‌شود.
\end{articlebox}

\chapter{مجلس مهستان (نهاد عالی گذار)}

\begin{articlebox}{ماده ۶: ترکیب و نمایندگی}
مجلس مهستان با ۳۰۰ کرسی تشکیل می‌شود:
۱. ۲۴۰ کرسی استانی (نسبی-فهرستی).
۲. ۴۰ کرسی فهرست‌های ملی سراسری.
۳. ۲۰ کرسی تضمینی اقلیت‌های قومی و مذهبی.
۴. سهمیه جنسیتی: حداقل ۳۰٪ فهرست‌ها باید به زنان اختصاص یابد (قاعده ۱ به ۳).
\end{articlebox}

\begin{articlebox}{ماده ۷: مشارکت ایرانیان خارج از کشور (دیاسپورا)}
ایرانیان خارج از کشور حق رأی مستقیم به فهرست‌های منتخب در حوزه‌های انتخاباتی داخلی را دارا می‌باشند. مجلس مهستان موظف به تسهیل زیرساخت‌های الکترونیک برای مشارکت آن‌هاست.
\end{articlebox}

\begin{articlebox}{ماده ۸: کمیسیون‌های تخصصی دائمی}
مجلس موظف به تشکیل ۱۱ کمسیون تخصصی از جمله: قانون اساسی موقت، عدالت انتقالی، اصلاح امنیتی (SSR)، اقتصاد و بودجه، حقوق زنان، و محیط‌زیست است.
\end{articlebox}

\chapter{ارکان اجرایی و قضایی انتقالی}

\begin{articlebox}{ماده ۹: کابینه اجرایی موقت}
۱. تشکیل کابینه بر عهده جریان اکثریت یا ائتلاف غالب است. در صورت بن‌بست، مجلس برای هر وزارتخانه به‌صورت فردی رأی‌گیری می‌کند.
۲. کابینه برای استیضاح به ۷۵٪ آرای مجلس نیاز دارد.
۳. پاسخگویی و شفافیت مالی تمامی اعضای کابینه الزامی است.
\end{articlebox}

\begin{articlebox}{ماده ۱۰: شورای قضایی انتقالی}
مدیریت دستگاه قضا به شورایی از قاضیان و حقوقدانان مستقل سپرده می‌شود. اولویت‌ها: تصفیه قضات متخلف، تعلیق قوانین ظالمانه و نظارت بر دادگاه‌های ویژه عدالت انتقالی.
\end{articlebox}

\begin{articlebox}{ماده ۱۱: هیأت میانجی‌گری مستقل}
جهت حل اختلافات احتمالی بین مجلس، کابینه و شورای قضایی، هیأتی متشکل از ۵ شخصیت مورد احترام ملی انتخاب می‌شود. در صورت عدم حل اختلاف، موضوع مستقیماً به رفراندوم گذاشته می‌شود.
\end{articlebox}

\chapter{عدالت انتقالی و اصلاح بخش امنیتی}

\begin{articlebox}{ماده ۱۲: ارکان پنج‌گانه‌ی عدالت انتقالی}
مجلس موظف به اجرای مدل جامع عدالت شامل این ۵ ستون است: حقیقت‌یابی (کمیسیون حقیقت)، عدالت کیفری (محاکمه آمران)، جبران خسارت (غرامت)، آشتی ملی (گفتگوی اجتماعی) و تضمین عدم تکرار (اصلاح نهادی).
\end{articlebox}

\begin{articlebox}{ماده ۱۳: اصلاح بخش امنیتی و دفاعی (SSR)}
۱. فرماندهی کل تحت نظارت مدنی مجلس است.
۲. انحلال قطعی نهادهای موازی و سرکوب‌گر و ادغام نیروهای غیرمتخلف در ارتش واحد ملی پس از فرایند غربال‌گری (Vetting).
۳. انتقال حسابرسی شده تمامی بنگاه‌های اقتصادی نظامی به صندوق ملی بازسازی.
\end{articlebox}

\chapter{نهادهای ناظر و تخصصی مستقل}

\begin{articlebox}{ماده ۱۴: آمبودزمان (نهاد حقوق شهروندی)}
نهادی مستقل برای رسیدگی به شکایات شهروندان از سوءمدیریت یا نقض حقوق در نهادهای دوران گذار که ریاست آن توسط مجلس انتخاب می‌شود.
\end{articlebox}

\begin{articlebox}{ماده ۱۵: کمیسیون ملی رسانه‌ها}
مسئول تبدیل صداوسیما به رسانه ملی مستقل و تضمین آزادی مطبوعات و جلوگیری از انحصار در فضای دیجیتال.
\end{articlebox}

\begin{articlebox}{ماده ۱۶: شورای اقتصادی دوران گذار}
مرکزی متشکل از متخصصان تراز اول برای مدیریت بحران، مهار تورم و شفاف‌سازی بودجه‌های کلان و حسابرسی بنیادهای غیردولتی.
\end{articlebox}

\chapter{فرایند تقنینی (ترکیب علم و دموکراسی)}

\begin{articlebox}{ماده ۱۷: بازوهای دوقلوی تقنینی}
تصویب هر قانون مشروط به دریافت دو گزارش موازی است:
۱. گزارش فنی از «مرکز تحقیقات» (بُعد واقعیت).
۲. گزارش بازخورد از «واحد ارتباط مردمی» (بُعد ارزش/نظرسنجی).
\end{articlebox}

\begin{articlebox}{ماده ۱۸: مراحل تصویب قانون}
هر طرح یا لایحه طی ۸ مرحله شامل ارجاع به کمیسیون، انتشار عمومی برای بازخورد (۷ تا ۳۰ روز)، استماع متخصصان، و رأی‌گیری در صحن (اکثریت ساده برای قوانین عادی و ۲/۳ برای قوانین مهم) نهایی می‌شود.
\end{articlebox}

\chapter{اقدامات اضطراری و اولویت‌های تقنینی (۱۰۰ روز اول)}

\begin{articlebox}{ماده ۱۹: فرمان‌های فوری هفته اول}
مجلس موظف به صدور ۱۰ فرمان: آزادی تمامی زندانیان سیاسی، لغو حجاب اجباری، رفع فیلترینگ، تعلیق اعدام، لغو دادگاه‌های انقلاب، حق بازگشت تبعیدیان، و حفاظت از اسناد امنیتی است.
\end{articlebox}

\begin{articlebox}{ماده ۲۰: اولویت‌های هفته ۱ تا ۴}
تصویب قوانین آزادی مطبوعات، آزادی احزاب و تظاهرات، لغو قوانین تبعیض‌آمیز علیه زنان، و حقوق فرهنگی اقلیت‌ها.
\end{articlebox}

\begin{articlebox}{ماده ۲۱: اولویت‌های ماه دوم تا چهارم}
تصویب قانون عدالت انتقالی، تشکیل کمیسیون حقیقت‌یابی، استقلال بانک مرکزی، و حسابرسی ملی بنیادها.
\end{articlebox}

\chapter{پایان مأموریت و انتقال قدرت}

\begin{articlebox}{ماده ۲۲: تدوین قانون اساسی نهایی}
مجلس مهستان موظف به تشکیل مجلس مؤسسان یا کمیسیون ویژه برای تدوین قانون اساسی نهایی با مشاوره‌ی گسترده عمومی است. متن نهایی الزاماً باید به رفراندوم ملی گذاشته شود.
\end{articlebox}

\begin{articlebox}{ماده ۲۳: حق عزل نماینده بدعهد (Recall)}
چنانچه ۴۰٪ از رأی‌دهندگان دور قبل در یک حوزه‌ی انتخاباتی خواستار عزل نماینده باشند، وی عزل و نفر بعدی لیست جایگزین می‌گردد.
\end{articlebox}

\begin{articlebox}{ماده ۲۴: انحلال و انتقال رسمی قدرت}
به محض تأیید نتایج انتخابات بر اساس قانون اساسی نهایی، مجلس مهستان منحل و تمامی اختیارات به نهادهای ناشی از قانون اساسی جدید منتقل می‌گردد.
\end{articlebox}

\end{document}
